\begin{abstract}
Vision-language models (VLMs) are often given an image together with text that may be wrong, such as a caption, a retrieved passage or a user's claim. We ask whether a model notices such a conflict, answers from the right input, and can say what is wrong, and whether these abilities go together. We introduce \name, a benchmark of 2,001 image--text pairs built from GQA scene graphs. In each conflict exactly one object, attribute or spatial relation in a faithful description is changed, graded into three levels of perceptual difficulty, and faithful controls measure false alarms. Every pair is probed for detection, resolution and attribution, with control probes that separate perception from arbitration and conversational structure from self-report. \pending{Two annotators validate 500 pairs.} Across eleven open models from four families, 2B to 32B parameters, the three abilities come apart. Models flag 72--98\% of conflicts, yet ``noticed, but not answered from the image'' is the most common outcome for six of them. Even when a model reads the fact correctly from the image alone, the best one answers from the image on only 73\% of conflicts. Scale helps up to 7--8B, but both 32B models answer from the image less often than their 7--8B siblings, while following an explicit instruction to rely on the image on 80--86\% of conflicts. Larger models also reject more faithful descriptions: false alarms rise from the 2--3B to the 7--8B model in every family, and within Qwen3-VL false-alarm-corrected detection falls from 59\% at 2B to 44\% at 32B. Asking a model to report the error before answering helps, but for nine of eleven models the turn structure alone contributes more than the report itself, and several models' error reports contradict their own detection answers. We will release \name, the code and all model outputs.
\end{abstract}
